\chapter{Fazit und Ausblick}
\label{ch:fazit}

\section{Fazit}
% Kurze Zusammenfassung: welche der 5 Fragestellungen wurden wie
% beantwortet, Gesamteinschaetzung.

Im Rahmen von WikiPod wurde untersucht, inwieweit sich ein lokaler Zugriff auf
ausgewählte Inhalte der englischsprachigen Wikipedia mit semantischem Retrieval
und lokaler Antwortgenerierung auf eingeschränkter Hardware umsetzen lässt.
Ausgehend von den fünf in der Aufgabenstellung formulierten Anforderungen wurde
dafür eine durchgängige Pipeline von der lokalen Wikipedia-Kopie über die
Dokumentenauswahl und Vektorindexierung bis hin zur Retrieval-Augmented-
Generation und deren Evaluation entwickelt.

Die erste Anforderung bestand in der Bereitstellung einer lokalen Kopie der
englischsprachigen Wikipedia. Hierfür werden Wikipedia-Inhalte in Form einer
lokal gespeicherten ZIM-Datei verwendet, die ohne bestehende Internetverbindung
verarbeitet werden kann. Damit bildet die ZIM-Datei die Grundlage für den
vollständig lokalen Zugriff auf den verfügbaren Wikipedia-Bestand.

Für die zweite Anforderung, die Entwicklung einer Methode zur Auswahl einer
Teilmenge der Wikipedia, wurde ein budgetbeschränktes Auswahlverfahren
implementiert. Die Artikel werden anhand mehrerer Metadaten und
Struktursignale bewertet und anschließend nach ihrem Verhältnis von Score zu
Speicherbedarf priorisiert. Im produktiven Selektionslauf auf dem vollständigen
Korpus wurden aus 7.322.195 Kandidatenartikeln 1.638.375 Artikel ausgewählt und
das konfigurierte Speicherbudget von 13.000~MB praktisch vollständig
ausgeschöpft. Die Analyse des Scoring-Modells zeigte jedoch zugleich, dass die
tatsächlichen Beiträge der einzelnen Signale nicht proportional zu ihren
konfigurierten Gewichten ausfallen. Insbesondere das Signal
\texttt{importance} besitzt einen deutlich größeren Einfluss auf den
Gesamtscore als ursprünglich durch seine Gewichtung beabsichtigt.

Die dritte Anforderung, die Indexierung der ausgewählten Teilmenge als dichte
Vektoren mittels OpenSearch, wurde durch eine Chunking- und Embedding-Pipeline
umgesetzt. Die ausgewählten Artikel werden abschnittsweise in Text-Chunks
zerlegt und mit \texttt{all-MiniLM-L6-v2} in dichte Vektorrepräsentationen
überführt \cite{minilm-model-card,sbert-pretrained-models}. Dense
Passage Retrieval zeigt grundsätzlich, dass dichte Repräsentationen für
semantisches Passage Retrieval im Open-Domain-Question-Answering eingesetzt
werden können \cite{karpukhin-etal-2020-dense}. Diese werden gemeinsam mit den
zugehörigen Metadaten in OpenSearch gespeichert. Für eingehende Nutzeranfragen
wird dasselbe Embedding-Modell verwendet, sodass über die Vektorsuche
semantisch ähnliche Textpassagen ermittelt und für die weitere Verarbeitung
bereitgestellt werden können.

Zur Beantwortung der vierten Anforderung wurden drei lokal ausführbare
Sprachmodelle auf dem Raspberry Pi hinsichtlich ihrer Generierungslatenz
verglichen. Qwen2.5~1.5B~Q4 erreichte mit durchschnittlich $27{,}97$~s die
niedrigste gemessene Generierungszeit, gefolgt von Llama~3.2~1B~Q4 mit
$29{,}18$~s und Phi-3.5~Mini~Q4 mit $152{,}33$~s. Für die Zielkonfiguration
wird daher Qwen2.5~1.5B in quantisierter Form über das
\texttt{llama\_cpp}-Backend eingesetzt. Diese Auswahl stützt sich auf die
gemessene Laufzeit und die erfolgreiche Ausführung auf der Zielhardware. Da
für den Modellvergleich keine quantitative Bewertung der Antwortqualität
durchgeführt wurde, kann daraus jedoch nicht abgeleitet werden, dass Qwen2.5
auch hinsichtlich der generierten Inhalte das qualitativ beste der
untersuchten Modelle ist.

Die fünfte Anforderung bestand in der Evaluation der entwickelten Lösung
hinsichtlich Effektivität und Effizienz. Auf dem regulären
Evaluationsdatensatz mit 20 Anfragen erreichte das Retrieval sowohl bei
$k=5$ als auch bei $k=10$ einen Recall von $1{,}000$, einen MRR von
$0{,}950$ und einen mittleren nDCG-Wert von $0{,}963$. Auf dem separaten
Edge-Case-Datensatz verbesserte der Query Analyzer den Recall@5 von
$0{,}833$ auf $0{,}917$ und den MRR von $0{,}792$ auf $0{,}875$. Die
mittlere Retrieval-Latenz auf dem Raspberry Pi betrug im untersuchten
Evaluationsaufbau $62{,}78$~ms und lag damit deutlich unter den
Generierungszeiten der lokalen Sprachmodelle. Die Antwortgenerierung stellt
somit in diesem Versuchsaufbau den zeitlich dominierenden Teil der
Anfragebearbeitung dar. Gleichzeitig ist die Aussagekraft der
Retrieval-Evaluation durch die geringe Größe und die manuelle Zusammenstellung
der verwendeten Evaluationsdatensätze begrenzt.

Zusammenfassend wurden alle fünf zentralen Anforderungen der Projektarbeit
umgesetzt beziehungsweise experimentell untersucht. Wikipedia-Inhalte können
lokal bereitgestellt, unter einem vorgegebenen Speicherbudget auf eine
Teilmenge reduziert, als dichte Vektoren in OpenSearch indexiert und für
semantisches Retrieval mit anschließender lokaler Antwortgenerierung genutzt
werden. Die Messungen zeigen, dass der grundsätzliche Betrieb der entwickelten
Pipeline auf dem Raspberry Pi möglich ist. Gleichzeitig bleiben insbesondere
die Gewichtung des Auswahlverfahrens, die begrenzte Breite der Evaluation und
die vergleichsweise hohe Laufzeit der lokalen Sprachmodellgenerierung als
wesentliche Einschränkungen der aktuellen Lösung bestehen.

\section{Ausblick}
% Was waere der naechste Schritt mit mehr Zeit -- z. B. voller Korpus auf
% dem Pi (Build/Serve-Trennung produktiv nutzen), Kategorien-basierte
% diversitaetsbewusste Selection, generierungsseitige Evaluationsmetriken
% (ROUGE/LLM-as-Judge) zusaetzlich zu Recall@k/MRR.

Ein naheliegender nächster Schritt besteht darin, die entwickelte Pipeline auf
einem deutlich größeren Wikipedia-Bestand unter der vorgesehenen
Zielkonfiguration zu betreiben. Die Architektur trennt die aufwendige
Aufbereitung und Indexierung von der späteren Nutzung des Systems, sodass ein
umfangreicher Index auf leistungsfähigerer Hardware erzeugt und anschließend
auf den Raspberry Pi übertragen werden kann. Dadurch ließe sich untersuchen,
wie sich die im kleinen Evaluationskorpus beobachteten Ergebnisse auf einen
realistischeren Datenbestand übertragen und wie sich insbesondere Indexgröße,
Speicherbedarf und Retrieval-Latenz bei wachsendem Korpus entwickeln.
Weiteres Verbesserungspotenzial besteht bei der Dokumentenauswahl. Die
Evaluation des Scoring-Verfahrens hat gezeigt, dass die tatsächlichen Beiträge
der verwendeten Signale nicht unmittelbar ihren konfigurierten Gewichten
entsprechen. Zukünftige Arbeiten sollten deshalb die Eingangssignale stärker
normalisieren und die Gewichte systematisch auf einem geeigneten
Evaluationsdatensatz bestimmen. Zusätzlich könnte eine kategorien- oder
themenbasierte Auswahl eingesetzt werden, um neben der geschätzten Relevanz
eines einzelnen Artikels auch die inhaltliche Diversität der ausgewählten
Wikipedia-Teilmenge zu berücksichtigen.
Auch die Retrieval-Evaluation sollte auf einen größeren und systematischer
zusammengestellten Datensatz erweitert werden. Neben zusätzlichen regulären
Anfragen sollten insbesondere mehrdeutige Entitäten, alternative
Bezeichnungen, Tippfehler, sehr kurze Anfragen und Out-of-Scope-Fälle stärker
vertreten sein. Für Out-of-Scope-Anfragen wäre zudem eine separate
Bewertungslogik sinnvoll, da die bisher verwendeten klassischen
Retrieval-Metriken Fälle ohne relevanten Ground-Truth-Artikel nicht angemessen
von gewöhnlichen Retrieval-Fehlern unterscheiden. Der Query Analyzer könnte
auf dieser Grundlage gezielter weiterentwickelt und seine Wirkung auf
unterschiedliche Fehlerklassen untersucht werden.
Schließlich sollte neben dem Retrieval auch die Qualität der generierten
Antworten systematisch evaluiert werden. Der bisher erprobte LLM-Judge zeigte
in einzelnen Fällen inkonsistente Bewertungen und eignet sich in seiner
aktuellen Form daher nicht als alleinige quantitative Qualitätsmetrik. Eine
zukünftige Evaluation könnte manuell bewertete Referenzantworten mit
automatisierten Verfahren kombinieren und dabei insbesondere Faktentreue,
Relevanz, Vollständigkeit und die korrekte Nutzung des bereitgestellten
Retrieval-Kontexts betrachten. Auf dieser Grundlage ließen sich die
untersuchten lokalen Sprachmodelle nicht nur hinsichtlich ihrer Laufzeit,
sondern auch hinsichtlich ihrer Antwortqualität vergleichen. Damit wäre eine
fundiertere Abwägung zwischen Ressourcenbedarf, Generierungslatenz und
Antwortqualität für den Einsatz auf dem Raspberry Pi möglich.
